\textbf{Stage~2 interleaving makes a two-source merge that no Stage~1 instance contains.} Rows i0--i5 are an STT Alloy instance, in the notation of Table~X: protected $m0$ (lavender) flows through i3 and the ALU op i4 into the address of transmitter i5 (thick arrows), and i3--i5 are speculative (tan band) because branch i2 is unresolved. i4's second input is left unspecified (?). Concretization step~4 (interleave) injects a chain of synthetic instructions, here c0 (a load) and c1 (an ALU op), at random positions after the committed prefix; chain steps carry no flags, and their outputs $f_0$, $f_1$ take registers from a free pool. Step~5 (register assignment) fills every unspecified operand from the same pool, so i4's ? can receive $f_1$'s register, which creates the dashed rf edge: i4 now merges c0's value, which is speculative only under the older branches i0 and i1, with i3's value, which stays speculative under i2. This is the shape behind the ReCon \texttt{getOldestTaint} bug: ReCon keeps only the oldest taint source of an ALU result, so once c0's speculation resolves it untaints i4's result even though i3's value ($r_0$) should still be protected, and i5 transmits.
